\documentclass[../../main.tex]{subfiles}
\begin{document}

\chapter{The Transformer architecture and attention mechanisms}
\label{ch:transformer_architecture}

\section{Chapter overview and learning outcomes}
In 2017, Vaswani et al. \citep{vaswani2017} introduced the Transformer in their seminal work \emph{"Attention Is All You Need"}, entirely discarding recurrence and convolutions in favor of self-attention mechanisms. This chapter presents an exhaustive architectural and mathematical breakdown of the Transformer model.

Upon completing this chapter, you will be able to:
\begin{itemize}
    \item Explain the fundamental bottleneck of sequential recurrence and how self-attention achieves $O(1)$ sequential operations during training.
    \item Derive Scaled Dot-Product Attention and explain the critical role of the $\sqrt{d_k}$ scaling factor.
    \item Formulate Multi-Head Attention (MHA) and contrast it with Multi-Query (MQA) and Grouped-Query Attention (GQA).
    \item Analyze Positional Encodings: absolute sinusoidal encodings, learned embeddings, and Rotary Position Embedding (RoPE).
    \item Dissect the sublayer structures: Residual connections, Post-LN vs Pre-LN, RMSNorm, and Feed-Forward Networks (FFN / SwiGLU).
    \item Compute the computational and memory complexity of self-attention ($O(T^2)$ vs $O(T)$).
\end{itemize}

\section{Scaled dot-product attention}
Let $\mQ \in \R^{T \times d_k}$, $\mK \in \R^{S \times d_k}$, and $\mV \in \R^{S \times d_v}$ denote Query, Key, and Value matrices.

\begin{definition}{Scaled dot-product attention}{scaled_dot_product}
The attention function maps queries and key-value pairs to an output matrix via:
\begin{equation}
\label{eq:scaled_dot_product}
\Attn{\mQ, \mK, \mV} = \softmax\left( \frac{\mQ \mK\trans}{\sqrt{d_k}} + \mat{M} \right) \mV,
\end{equation}
where $\mat{M} \in \{0, -\infty\}^{T \times S}$ is an optional causal mask matrix.
\end{definition}

\begin{examinsight}[Exam highlight: why scale by $\sqrt{d_k}$?]
A quintessential oral exam question asks: \emph{"What is the exact mathematical justification for dividing $\mQ \mK\trans$ by $\sqrt{d_k}$?"}
Assume components of $\vq$ and $\vk$ are independent random variables with mean $0$ and variance $1$. The dot product $\vq\trans \vk = \sum_{i=1}^{d_k} q_i k_i$ has mean $0$ and variance $\sum_{i=1}^{d_k} 1 = d_k$.
For large $d_k$ (e.g., $d_k = 64$ or $128$), variance is large, pushing the dot products into regions where the $\softmax$ function has extremely small gradients (saturation). Dividing by $\sqrt{d_k}$ normalizes the variance back to $1$, preserving healthy gradient flow during backpropagation.
\end{examinsight}

\begin{intuition}[Attention as differentiable soft key-value retrieval]
In computer systems, a hash table maps a discrete key to a unique value. Self-attention is a soft, differentiable generalisation: a query $\vq$ is compared against all keys $\vk_i$ via inner products. The resulting softmax weights dictate how much of each value vector $\vv_i$ is blended into the output representation.
\end{intuition}

\section{Multi-head attention and sublayer design}
\begin{definition}{Multi-head attention (MHA)}{mha}
Instead of performing a single attention function with $d_{\text{model}}$-dimensional queries, keys, and values, Multi-Head Attention projects them linearly $h$ times with different learned parameter matrices:
\begin{equation}
\MHA{\mQ, \mK, \mV} = \left[ \text{head}_1 \parallel \text{head}_2 \parallel \dots \parallel \text{head}_h \right] \mW^O,
\end{equation}
where $\text{head}_i = \Attn{\mQ \mW_i^Q, \mK \mW_i^K, \mV \mW_i^V}$, with $\mW_i^Q \in \R^{\dmodel \times \dk}$, $\mW_i^K \in \R^{\dmodel \times \dk}$, $\mW_i^V \in \R^{\dmodel \times \dv}$, and $\mW^O \in \R^{h \dv \times \dmodel}$.
\end{definition}

\begin{deepdive}[Vaswani et al. (2017) and modern Transformer variants]
The original Transformer architecture \citep{vaswani2017} employed Post-LayerNorm and sinusoidal positional encodings. Modern LLM architectures (such as LLaMA, Mistral, and DeepSeek) have evolved to utilize Pre-LayerNorm with RMSNorm, Rotary Position Embeddings (RoPE), and SwiGLU activations, improving training stability across hundreds of billions of tokens.
\end{deepdive}

\end{document}
